\chapter{The Linear Algebra Method}

\begin{center}\itshape
It is my experience that proofs involving matrices can be shortened by 50\% if one throws the matrices out.\\
\upshape(Emil Artin, \emph{Geometric Algebra}, 1957)
\end{center}

A town has $n$ inhabitants and a collection of clubs. Every club contains an odd number of people, while any two different clubs overlap in an even number of people. How many clubs can there be?

Give each club a vector in $\mathbb F_2^n$: write $1$ for a member and $0$ for a non-member. The dot product of a club's vector with itself is $1$, because the club has odd size. The dot product of two different clubs is $0$, because they overlap in an even number of places.

Suppose a linear combination of these club vectors were zero. Taking the dot product with any one vector shows that its coefficient must be zero. Every coefficient is zero, so the vectors are linearly independent. There can be no more than $n$ clubs.

This proof uses a dimension bound where a direct counting argument would struggle. The rest of the appendix develops that trick with polynomials, intersection patterns and finite fields. We will keep asking the same question: what vector space is too small to hold all the objects we want?

\section{Oddtown, Eventown, and Fisher}

For $A\subseteq[n]=\{1,\dots,n\}$ write $v_A\in\{0,1\}^n$ for its \emph{characteristic vector}, with a $1$ in position $i$ when $i\in A$. The key identity of the
whole chapter is
\[
v_A\cdot v_B=\sum_{i=1}^n(v_A)_i(v_B)_i=|A\cap B|,
\]
which we may read in $\R$ or, reducing modulo a prime $p$, in the field $\mathbb F_p=\Z/p\Z$.

\begin{theorem}[Oddtown, Berlekamp, 1969]\label{combinatorial:thm:oddtown}
Let $A_1,\dots,A_m\subseteq[n]$ be such that every $|A_i|$ is odd and every $|A_i\cap A_j|$, $i\ne j$, is even. Then $m\le n$.
\end{theorem}
\begin{proof}
Work in $\mathbb F_2^n$ and write $v_i=v_{A_i}$. The hypotheses say $v_i\cdot v_j=1$ if $i=j$ and $0$ otherwise. If $\sum_ic_iv_i=0$ with $c_i\in\mathbb F_2$, taking
the dot product with $v_j$ gives $c_j=0$. So $v_1,\dots,v_m$ are linearly independent in a space of dimension $n$, and $m\le n$.
\end{proof}

Eventown is different. Split the inhabitants into couples, $\{1,2\},\{3,4\},\dots$, and let the clubs be all unions of couples: every club and every
intersection of two clubs is a union of couples, of even size, and there are $2^{\lfloor n/2\rfloor}$ clubs. Changing one word in the rules has
changed the answer from $n$ to an exponential. Linear algebra sees the difference too: in Eventown, the vectors are orthogonal \emph{to themselves}.

\begin{lemma}\label{combinatorial:lem:perp}
Let $F$ be a field and $W$ a subspace of $F^n$. Then $W^\perp=\{u\in F^n:u\cdot w=0\text{ for all }w\in W\}$ is a subspace of dimension $n-\dim W$.
\end{lemma}
\begin{proof}
Let $w_1,\dots,w_d$ be a basis of $W$ and $M$ the $d\times n$ matrix with rows $w_i$. Then $W^\perp$ is the kernel of $u\mapsto Mu$. The rank of $M$ is $d$ (its rows are
independent, and row rank equals column rank), so by the rank--nullity theorem its kernel has dimension $n-d$.
\end{proof}

Over $\mathbb F_2$ (or $\mathbb C$, or any field other than $\R$) a non-zero vector can be orthogonal to itself, so $W$ and $W^\perp$ may intersect, and even
$W\subseteq W^\perp$ is possible: $(1,1)\cdot(1,1)=0$ in $\mathbb F_2^2$. The lemma still bounds how large such a $W$ can be.

\begin{theorem}[Eventown, Berlekamp, 1969; Graver, 1975]\label{combinatorial:thm:eventown}
Let $A_1,\dots,A_m\subseteq[n]$ be distinct sets, all of even size, with every $|A_i\cap A_j|$ even. Then $m\le2^{\lfloor n/2\rfloor}$.
\end{theorem}
\begin{proof}
Let $W$ be the span of $v_{A_1},\dots,v_{A_m}$ in $\mathbb F_2^n$. All the products $v_{A_i}\cdot v_{A_j}$ vanish, including for $i=j$, so by bilinearity $w\cdot w'=0$ for
all $w,w'\in W$: $W\subseteq W^\perp$. By Lemma~\ref{combinatorial:lem:perp}, $\dim W\le n-\dim W$, so $\dim W\le\lfloor n/2\rfloor$. The $m$ distinct vectors lie in $W$, which has
$2^{\dim W}$ elements.
\end{proof}

Both bounds are attained, by singletons and by couples, but not only by them: for $n=8$, the $14$ affine planes of $\mathbb F_2^3$, with $\varnothing$ and the whole space, form an Eventown family of $16=2^4$ clubs
that contains no couple at all (Exercise~\ref{combinatorial:ex:eventown}). It is a famous code in disguise, the extended Hamming code, and the self-orthogonal spaces of Theorem~\ref{combinatorial:thm:eventown} are the
\emph{self-dual codes} of coding theory.

Over the reals, the method gives one of the founding results of the design of experiments. Ronald Fisher spent fourteen years, from 1919 to 1933, designing field trials at the Rothamsted
agricultural station. In 1940 he studied plans of the following kind: $v$ varieties of a crop are compared on a field split into blocks, each block can hold only some of the varieties, and for a fair comparison any two varieties
should appear together in the same number of blocks. Fisher proved that such a plan needs at least as many blocks as varieties. In the language of sets (exchanging
the roles of varieties and blocks, by counting incidences), this is the following theorem.

\begin{theorem}[Fisher's inequality, non-uniform form; Majumdar, 1953]\label{combinatorial:thm:fisher}
Let $\lambda\ge1$, and let $A_1,\dots,A_m$ be distinct subsets of $[n]$ such that $|A_i\cap A_j|=\lambda$ for all $i\ne j$. Then $m\le n$.
\end{theorem}
\begin{proof}
Let $M$ be the $m\times n$ real matrix with rows $v_{A_i}$. Its Gram matrix $G=MM^{\mathsf T}$ has entries $|A_i\cap A_j|$, so
\[
G=\lambda J+D,\qquad D=\mathrm{diag}(d_1,\dots,d_m),\quad d_i=|A_i|-\lambda,
\]
where $J$ is the all-ones matrix. Each $d_i\ge0$, since $|A_i|\ge|A_i\cap A_j|=\lambda$ (if $m=1$ there is nothing to prove). At most one $d_i$ is zero: if $|A_i|=\lambda$
then $A_i\subseteq A_j$ for every $j$, and two such sets would contain each other, hence be equal. For $x\in\R^m$,
\[
x^{\mathsf T}Gx=\lambda\Bigl(\sum_ix_i\Bigr)^2+\sum_id_ix_i^2.
\]
If this vanishes, then $d_ix_i=0$ for all $i$, so $x$ is zero outside the at most one index $i_0$ with $d_{i_0}=0$, and then $\lambda x_{i_0}^2=0$ gives $x=0$. So $G$ is positive definite,
hence invertible, and $m=\operatorname{rank}G\le\operatorname{rank}M\le n$.
\end{proof}

\begin{think}
Where exactly did the proof use $\lambda\ge1$? Find a family with $\lambda=0$ and $m>n$.
\end{think}

\section{Polynomials are vectors too}

The second idea of the method is that functions form vector spaces. To bound the size of a set $\{a_1,\dots,a_m\}$, attach to each point $a_i$ a function $f_i$ that does not vanish at
$a_i$ but vanishes at all the other points. Then $f_1,\dots,f_m$ are linearly independent: if $\sum_ic_if_i=0$, evaluating at $a_j$ leaves $c_jf_j(a_j)=0$. If all the $f_i$ lie in a
space of polynomials of low degree, $m$ is at most its dimension. This is the polynomial version of a pairing argument, and it is remarkably flexible: it is enough that the matrix
$\bigl(f_i(a_j)\bigr)$ be \emph{triangular} with non-zero diagonal.

\subsection*{Few distances}

A \emph{two-distance set} in $\R^n$ is a set of points whose pairwise distances take only two values. The vertices of a regular pentagon are a two-distance set in the plane (Figure~\ref{combinatorial:fig:pentagon}).
How large can a two-distance set in $\R^n$ be?

\begin{figure}[ht]
\centering
\fig{combinatorial-ch11-pentagon}
\caption{The regular pentagon is a two-distance set: sides (charcoal) and diagonals (coral). In the plane no two-distance set has more than $5$ points (Kelly, 1947). Theorem~\ref{combinatorial:thm:twodist} gives $9$, Blokhuis' refinement $6$.}
\label{combinatorial:fig:pentagon}
\end{figure}

\begin{theorem}[Larman, Rogers and Seidel, 1977]\label{combinatorial:thm:twodist}
A two-distance set in $\R^n$ has at most $\frac{(n+1)(n+4)}2$ points.
\end{theorem}
\begin{proof}
Let $a_1,\dots,a_m$ have pairwise distances in $\{b,c\}$, with $b,c>0$. Define
\[
f_i(x)=\bigl(|x-a_i|^2-b^2\bigr)\bigl(|x-a_i|^2-c^2\bigr),\qquad x\in\R^n.
\]
Then $f_i(a_j)=0$ for $j\ne i$ and $f_i(a_i)=b^2c^2\ne0$, so $f_1,\dots,f_m$ are linearly independent. Now $|x-a|^2=|x|^2-2a\cdot x+|a|^2$ is $|x|^2$ plus an affine function of $x$, so
\[
f_i(x)=\bigl(|x|^2+\ell_1(x)\bigr)\bigl(|x|^2+\ell_2(x)\bigr)=|x|^4+|x|^2\bigl(\ell_1(x)+\ell_2(x)\bigr)+\ell_1(x)\ell_2(x),
\]
with $\ell_1,\ell_2$ affine. Every $f_i$ is therefore a linear combination of the functions
\[
|x|^4,\quad |x|^2x_k\ (1\le k\le n),\quad x_kx_l\ (1\le k\le l\le n),\quad x_k\ (1\le k\le n),\quad 1,
\]
where we used that $|x|^2=\sum_kx_k^2$ is itself a combination of the $x_kx_l$. There are $1+n+\frac{n(n+1)}2+n+1=\frac{(n+1)(n+4)}2$ of them.
\end{proof}

The truth is not far: the midpoints of the edges of a regular simplex give $\frac{n(n+1)}2$ points (Exercise~\ref{combinatorial:ex:simplexdist}), so the maximum grows like $\frac{n^2}2$. Blokhuis (1984) improved the bound to
$\frac{(n+1)(n+2)}2$ by showing that the functions $f_i$ remain independent even after adding the $n+1$ functions $x_1,\dots,x_n,1$; in dimensions $2$ and $3$ the maxima are $5$ and $6$ (the regular octahedron).
The same proof bounds $s$-distance sets (Exercise~\ref{combinatorial:ex:sdistance}).

\subsection*{Intersections modulo a prime}

When the points are vertices of the cube $\{0,1\}^n$, that is, sets, the polynomials can be simplified, because $x_k^2=x_k$ on the cube. Call a polynomial \emph{multilinear} if no variable appears with
exponent more than $1$. Replacing each $x_k^e$, $e\ge1$, by $x_k$ in a polynomial $f$ gives a multilinear polynomial $\tilde f$ with the same values on $\{0,1\}^n$ and with $\deg\tilde f\le\deg f$; and the
multilinear polynomials of degree at most $s$ are combinations of the $\prod_{k\in K}x_k$ with $|K|\le s$, of which there are $\sum_{i=0}^s\binom ni$.

\begin{theorem}[Frankl and Wilson, 1981, modular form]\label{combinatorial:thm:fwmod}
Let $p$ be a prime and $L$ a set of $s$ residues modulo $p$. Let $A_1,\dots,A_m\subseteq[n]$ be such that $|A_i|\bmod p\notin L$ for all $i$, and $|A_i\cap A_j|\bmod p\in L$ for all $i\ne j$. Then
\[
m\le\sum_{i=0}^s\binom ni.
\]
\end{theorem}
\begin{proof}
Over $\mathbb F_p$, let $f_i(x)=\prod_{l\in L}(v_{A_i}\cdot x-l)$, a polynomial of degree $s$ in $x\in\mathbb F_p^n$. At the vertex $v_{A_j}$ it takes the value $\prod_{l\in L}(|A_i\cap A_j|-l)$, which is $0$ for $j\ne i$ and non-zero for
$j=i$. So the multilinear reductions $\tilde f_1,\dots,\tilde f_m$, which have the same values on the cube, are linearly independent, and they lie in the space spanned by the $\sum_{i\le s}\binom ni$
multilinear monomials of degree at most $s$.
\end{proof}

Oddtown is the case $p=2$, $L=\{0\}$, with the bound $1+n$; a little more care gives back $n$. Without the prime, a weaker condition on the intersection sizes alone is enough, at the price of using a triangular
rather than a diagonal system.

\begin{theorem}[Frankl and Wilson, 1981, non-uniform form]\label{combinatorial:thm:fw}
Let $L$ be a set of $s$ non-negative integers and $A_1,\dots,A_m\subseteq[n]$ distinct sets with $|A_i\cap A_j|\in L$ for all $i\ne j$. Then $m\le\sum_{i=0}^s\binom ni$.
\end{theorem}
\begin{proof}
Number the sets so that $|A_1|\le|A_2|\le\dots\le|A_m|$, and over $\R$ let
\[
f_i(x)=\prod_{l\in L,\ l<|A_i|}\bigl(v_{A_i}\cdot x-l\bigr),
\]
of degree at most $s$. Then $f_i(v_{A_i})=\prod_{l<|A_i|}(|A_i|-l)>0$. If $j<i$, then $|A_j|\le|A_i|$ and $A_j\ne A_i$, so $A_i\not\subseteq A_j$ and $|A_i\cap A_j|<|A_i|$; as this number is in $L$, one factor of
$f_i(v_{A_j})$ vanishes. Suppose $\sum_ic_i\tilde f_i=0$ with not all $c_i$ zero, and let $i_0$ be the \emph{smallest} index with $c_{i_0}\ne0$. Evaluate at $v_{A_{i_0}}$: the terms $i<i_0$ have $c_i=0$, the terms $i>i_0$ have
$\tilde f_i(v_{A_{i_0}})=0$, and what remains is $c_{i_0}f_{i_0}(v_{A_{i_0}})\ne0$, a contradiction. So the $\tilde f_i$ are independent, and we conclude as before.
\end{proof}

\subsection*{Explicit Ramsey graphs}

In Volume~4, Chapter~12, Erd\H os' coin flips showed that some graph on $2^{t/2}$ vertices has no clique and no independent set of size $t$. Nobody knows how to \emph{write down} such a graph: Erd\H os offered
\$100 for an explicit construction with $(1+\varepsilon)^t$ vertices, and the problem is still open. The first construction that beats every polynomial in $t$ came from Theorem~\ref{combinatorial:thm:fwmod}.

\begin{theorem}[Frankl and Wilson, 1981]\label{combinatorial:thm:explicitramsey}
Let $p$ be a prime, and let $G_p$ be the graph whose vertices are the $(p^2-1)$-element subsets of $[p^3]$, two of them $A,B$ being adjacent when $|A\cap B|\not\equiv-1\pmod p$. Then $G_p$ has no clique and no
independent set with more than $t_p=\sum_{i=0}^{p-1}\binom{p^3}i$ vertices.
\end{theorem}
\begin{proof}
A clique is a family of sets with $|A|=p^2-1\equiv-1$ and $|A\cap B|\bmod p\in L=\{0,1,\dots,p-2\}$, which has $p-1$ elements and avoids $-1$: Theorem~\ref{combinatorial:thm:fwmod} applies with $s=p-1$. In an independent set, all
intersections satisfy $|A\cap B|\equiv-1\pmod p$ and $|A\cap B|<p^2-1$, so $|A\cap B|\in\{p-1,2p-1,\dots,p^2-p-1\}$, a set of $p-1$ integers: Theorem~\ref{combinatorial:thm:fw} applies with $s=p-1$.
\end{proof}

For $p=3$ the graph has $\binom{27}8=2\,220\,075$ vertices and neither cliques nor independent sets with more than $1+27+351=379$ vertices. In general $t_p\le p\cdot p^{3(p-1)}\le p^{3p}$ while the number of vertices is
$\binom{p^3}{p^2-1}\ge\bigl(\frac{p^3}{p^2-1}\bigr)^{p^2-1}\ge p^{p^2-1}$. Since $\frac{(p^2-1)\log p}{3p\log p}\to\infty$, the number of vertices eventually exceeds $t_p^c$ for every fixed $c$: the growth is like $t^{c\log t/\log\log t}$,
superpolynomial, though far below $2^{t/2}$. Written in terms of the number $N$ of vertices, the largest clique or independent set of $G_p$ has size about $e^{c\sqrt{\log N\log\log N}}$, while for a random graph it is about $2\log_2N$. Explicit constructions
improved slowly, through ideas from the theory of randomness extractors : Chattopadhyay and Zuckerman (2016) reached $e^{(\log\log N)^{C}}$, and Li (2023) a power of $\log N$.

\begin{remark}[Borsuk's conjecture]
In 1933 Karol Borsuk asked whether every bounded set in $\R^d$ can be cut into $d+1$ pieces of smaller diameter, as a ball can. Sixty years later Kahn and Kalai (1993) answered no, in dimension $1325$,
with a set of $\pm1$-vectors built from Theorem~\ref{combinatorial:thm:fwmod}: pieces of smaller diameter correspond to families avoiding one intersection size, which are small, so many pieces are needed. Counterexamples are
now known down to dimension $64$ (Jenrich and Brouwer, 2014); the conjecture is open between $4$ and $63$.
\end{remark}

\section{The combinatorial Nullstellensatz}

A non-zero polynomial of degree $d$ in one variable has at most $d$ roots, so it cannot vanish on a set of $d+1$ points. In 1999 Noga Alon turned the several-variable version of this triviality into a tool
of startling power.

\begin{lemma}\label{combinatorial:lem:grid}
Let $F$ be a field, $S_1,\dots,S_n\subseteq F$ finite, and $f\in F[x_1,\dots,x_n]$ with $\deg_{x_i}f<|S_i|$ for every $i$. If $f$ vanishes on $S_1\times\dots\times S_n$, then $f=0$.
\end{lemma}
\begin{proof}
Induction on $n$; for $n=1$, a non-zero $f$ would have more roots than its degree. For $n>1$ write $f=\sum_{j<|S_n|}g_j(x_1,\dots,x_{n-1})\,x_n^j$. For fixed $(s_1,\dots,s_{n-1})\in S_1\times\dots\times S_{n-1}$,
the one-variable polynomial $\sum_jg_j(s_1,\dots,s_{n-1})x_n^j$ vanishes on $S_n$ and has degree less than $|S_n|$, so all its coefficients vanish: each $g_j$ vanishes on $S_1\times\dots\times S_{n-1}$, and by induction $g_j=0$.
\end{proof}

\begin{theorem}[Combinatorial Nullstellensatz, Alon, 1999]\label{combinatorial:thm:cn}
Let $F$ be a field and $f\in F[x_1,\dots,x_n]$. Suppose that $\deg f=t_1+\dots+t_n$ and that the coefficient of $x_1^{t_1}\cdots x_n^{t_n}$ in $f$ is non-zero. If $S_1,\dots,S_n\subseteq F$ satisfy $|S_i|>t_i$, then
there is $s\in S_1\times\dots\times S_n$ with $f(s)\ne0$.
\end{theorem}
\begin{proof}
Shrinking the $S_i$, we may assume $|S_i|=t_i+1$. Let $g_i(y)=\prod_{a\in S_i}(y-a)$, of degree $t_i+1$. For $e\ge0$, let $r_{i,e}(y)$ be the remainder of $y^e$ in the division by $g_i$: it has degree at most $t_i$, it
equals $y^e$ if $e\le t_i$, it has degree less than $e$ if $e>t_i$, and $r_{i,e}(a)=a^e$ for $a\in S_i$, since $g_i(a)=0$. Replace each monomial $x_1^{e_1}\cdots x_n^{e_n}$ of $f$ by $r_{1,e_1}(x_1)\cdots r_{n,e_n}(x_n)$; this gives
a polynomial $\bar f$ with $\deg_{x_i}\bar f\le t_i$ that agrees with $f$ on $S_1\times\dots\times S_n$.

Look at the coefficient of $x_1^{t_1}\cdots x_n^{t_n}$ in $\bar f$. The monomial itself is unchanged by the replacement. Any other monomial $m$ of $f$ either has all $e_i\le t_i$, and is unchanged and different, or has
some $e_i>t_i$, and is replaced by a polynomial of degree less than $\deg m\le\deg f=t_1+\dots+t_n$, which cannot contain $x_1^{t_1}\cdots x_n^{t_n}$. So $\bar f$ keeps the non-zero coefficient, $\bar f\ne0$, and by
Lemma~\ref{combinatorial:lem:grid}, $\bar f$, hence $f$, does not vanish everywhere on $S_1\times\dots\times S_n$.
\end{proof}

The hypothesis that $\deg f$ equals $\sum t_i$ matters: it guarantees that the chosen monomial is of top degree, so lower-order terms cannot cancel it. Here is the first application, a theorem of 1813.

\begin{theorem}[Cauchy, 1813; Davenport, 1935]\label{combinatorial:thm:cd}
Let $p$ be a prime and $A,B$ non-empty subsets of $\Z/p\Z$. Then $|A+B|\ge\min(p,|A|+|B|-1)$, where $A+B=\{a+b:a\in A,b\in B\}$.
\end{theorem}
\begin{proof}
If $|A|+|B|>p$, then for every $c$ the sets $A$ and $c-B$ have more than $p$ elements in total, so they meet: $c=a+b$, and $A+B=\Z/p\Z$. Suppose now $|A|+|B|\le p$ and, for a contradiction,
$|A+B|\le|A|+|B|-2$. Choose $C\supseteq A+B$ with $|C|=|A|+|B|-2$ and let
\[
f(x,y)=\prod_{c\in C}(x+y-c)\in\mathbb F_p[x,y],
\]
of degree $|A|+|B|-2$. The coefficient of $x^{|A|-1}y^{|B|-1}$ in $f$ is that of the top part $(x+y)^{|A|+|B|-2}$, the binomial coefficient $\binom{|A|+|B|-2}{|A|-1}$, which is not divisible by $p$ because
$|A|+|B|-2<p$. By Theorem~\ref{combinatorial:thm:cn} with $t_1=|A|-1$, $t_2=|B|-1$, there are $a\in A$, $b\in B$ with $f(a,b)\ne0$. But $a+b\in C$, so $f(a,b)=0$.
\end{proof}

The bound is sharp (arithmetic progressions with the same difference), and it fails for $\Z/n\Z$ with $n$ composite (Exercise~\ref{combinatorial:ex:cdsharp}): primes are needed so that $\Z/p\Z$ has no subgroups to hide in.

\subsection*{Counting zeros modulo $p$}

\begin{lemma}\label{combinatorial:lem:powersum}
Let $p$ be a prime and $0\le e<p-1$. Then $\sum_{y\in\mathbb F_p}y^e=0$ in $\mathbb F_p$ (with the convention $0^0=1$).
\end{lemma}
\begin{proof}
For $e=0$ the sum is $p\cdot1=0$. For $1\le e<p-1$, the polynomial $y^e-1$ has at most $e<p-1$ roots, so some $a\ne0$ has $a^e\ne1$. Multiplication by $a$ permutes $\mathbb F_p$, so
$S=\sum_yy^e=\sum_y(ay)^e=a^eS$, and $(a^e-1)S=0$ forces $S=0$.
\end{proof}

\begin{theorem}[Chevalley and Warning, 1935]\label{combinatorial:thm:cw}
Let $p$ be a prime and $P_1,\dots,P_k\in\mathbb F_p[x_1,\dots,x_n]$ with $\deg P_1+\dots+\deg P_k<n$. Then the number $N$ of common zeros of $P_1,\dots,P_k$ in $\mathbb F_p^n$ is divisible by $p$. In particular, if they
have the common zero $0$, they have another one.
\end{theorem}
\begin{proof}
By Fermat's little theorem, $y^{p-1}=1$ for $y\ne0$ in $\mathbb F_p$. So $Q=\prod_{j=1}^k\bigl(1-P_j^{p-1}\bigr)$ takes the value $1$ at common zeros and $0$ elsewhere, and $N\equiv\sum_{x\in\mathbb F_p^n}Q(x)\pmod p$. Now
$\deg Q\le(p-1)\sum_j\deg P_j<(p-1)n$, so every monomial $x_1^{e_1}\cdots x_n^{e_n}$ of $Q$ has some $e_i<p-1$, and
\[
\sum_{x\in\mathbb F_p^n}x_1^{e_1}\cdots x_n^{e_n}=\prod_{i=1}^n\Bigl(\sum_{y\in\mathbb F_p}y^{e_i}\Bigr)=0
\]
by Lemma~\ref{combinatorial:lem:powersum}. Hence $N\equiv0\pmod p$; if $N\ge1$ then $N\ge p\ge2$.
\end{proof}

The theorem holds over every finite field, with the same proof. It is a statement about solving equations, proved without solving anything; here is a statement about sums of integers, proved by solving equations.

\begin{theorem}[Erd\H os, Ginzburg and Ziv, 1961]\label{combinatorial:thm:egz}
Among any $2n-1$ integers there are $n$ whose sum is divisible by $n$.
\end{theorem}
\begin{proof}
First let $n=p$ be prime, and $a_1,\dots,a_{2p-1}$ the integers. Over $\mathbb F_p$, in $2p-1$ variables, consider
\[
P_1=\sum_{i=1}^{2p-1}a_ix_i^{p-1},\qquad P_2=\sum_{i=1}^{2p-1}x_i^{p-1}.
\]
Their degrees add up to $2p-2<2p-1$ and $0$ is a common zero, so by Theorem~\ref{combinatorial:thm:cw} there is a common zero $x\ne0$. Let $I=\{i:x_i\ne0\}$. By Fermat, $x_i^{p-1}$ is $1$ on $I$ and $0$ elsewhere, so
$P_2(x)=|I|\equiv0\pmod p$ with $1\le|I|\le2p-1$, that is $|I|=p$; and $P_1(x)=\sum_{i\in I}a_i\equiv0\pmod p$.

Next, if the theorem holds for $m$ and for $n$, it holds for $mn$. Given $2mn-1$ integers, as long as at least $2n-1$ remain, pick $n$ of them with sum divisible by $n$ and set them aside. After setting aside
$k$ groups, $2mn-1-kn$ remain, which is at least $2n-1$ while $k\le2m-2$: so we obtain $2m-1$ disjoint groups, with sums $nb_1,\dots,nb_{2m-1}$. Among the $b_j$, choose $m$ with sum divisible by $m$; the
$mn$ integers of the corresponding groups have a sum divisible by $mn$. Every $n\ge2$ is a product of primes, and $n=1$ is trivial.
\end{proof}

With $n-1$ zeros and $n-1$ ones, $2n-2$ integers are not enough (Exercise~\ref{combinatorial:ex:egzsharp}).

\section{Cap sets and the card game SET}

In the card game SET, invented in 1974 by the geneticist Marsha Falco while she was coding information about epileptic dogs, each of the $81$ cards shows one, two or three symbols, of one of three shapes,
three colours and three shadings. Three cards form a \emph{set} when, for each of the four attributes, they are all the same or all different. Number the values of each attribute $0,1,2$: a card is a vector
of $\mathbb F_3^4$, and three cards $x,y,z$ form a set exactly when $x+y+z=0$ in each coordinate (three equal values add to $0$, and $0+1+2=3=0$). Since $x+y+z=0$ means $z=-x-y$, two cards determine the
third: the sets are the \emph{lines} of the affine space $\mathbb F_3^4$, each containing $3$ points.

\begin{definition}\label{combinatorial:def:cap}
A \emph{cap set} in $\mathbb F_3^n$ is a set of points containing no line, that is, no three distinct points $x,y,z$ with $x+y+z=0$.
\end{definition}

Players often wonder how many cards can be on the table without any set. The answer is $20$ (Pellegrino, 1971); in $\mathbb F_3^1,\dots,\mathbb F_3^6$ the maximal sizes are $2,4,9,20,45,112$, the last two found by
heavy computation (Edel, Ferret, Landjev and Storme, 2002; Potechin, 2008). Figure~\ref{combinatorial:fig:caps} shows the case $n=2$, proved in Exercise~\ref{combinatorial:ex:capplane}. How do the maxima grow? Meshulam (1995) showed, by
Fourier analysis, that they are at most of order $\frac{3^n}n$: a vanishing proportion. The best constructions give about $2.22^n$. For twenty years the question was whether the true growth is $3^{n-o(n)}$ or
$c^n$ with $c<3$; many believed the former. In May 2016, Croot, Lev and Pach solved the analogous problem in $(\Z/4\Z)^n$ with polynomials, and within ten days Ellenberg and Gijswijt adapted the argument to
$\mathbb F_3^n$. Their papers are three and five pages long. We follow Terence Tao's symmetric reformulation, which isolates a notion of rank for functions of three variables.

\begin{figure}[ht]
\centering
\fig{combinatorial-ch11-caps}
\caption{The plane $\mathbb F_3^2$. Its $12$ lines come in four parallel classes of three lines (one colour per line in each panel; the last two classes are the ``diagonals'', which wrap around). Right: a cap of
size $4$, which meets every line in at most two points. No cap has $5$ points (Exercise~\ref{combinatorial:ex:capplane}).}
\label{combinatorial:fig:caps}
\end{figure}

\begin{definition}\label{combinatorial:def:slicerank}
Let $A$ be a finite set and $F$ a field. A function $T:A\times A\times A\to F$ is a \emph{slice} if it has one of the forms $f(x)g(y,z)$, $f(y)g(x,z)$ or $f(z)g(x,y)$. The \emph{slice rank} of $T$ is the least $r$ such
that $T$ is a sum of $r$ slices.
\end{definition}

For functions of two variables, the analogous notion (sums of products $f(x)g(y)$) is the rank of the matrix $(T(x,y))$. A diagonal matrix with non-zero diagonal has full rank; the key lemma says the same for
diagonal functions of three variables.

\begin{lemma}\label{combinatorial:lem:support}
A subspace $V\subseteq F^A$ of dimension $d$ contains a vector with at least $d$ non-zero coordinates.
\end{lemma}
\begin{proof}
Put a basis of $V$ as the rows of a $d\times|A|$ matrix in reduced row echelon form (Gaussian elimination keeps the row space). Each row has a $1$ in its own pivot column and $0$ in the other pivot
columns, so the sum of the rows, which is in $V$, has a $1$ in each of the $d$ pivot columns.
\end{proof}

\begin{lemma}[Tao, 2016]\label{combinatorial:lem:diag}
Let
\[
T(x,y,z)=\sum_{a\in A}c_a\,\delta_a(x)\delta_a(y)\delta_a(z),
\]
with all $c_a\ne0$, where $\delta_a(x)$ is $1$ if $x=a$ and $0$ otherwise. Then the slice rank of $T$ is $|A|$.
\end{lemma}
\begin{proof}
It is at most $|A|$, by the formula. Suppose $T=\sum_{i=1}^rS_i$ with $r<|A|$ slices: $r_1$ of the form $f_i(x)g_i(y,z)$, $r_2$ of the form $f_i(y)g_i(x,z)$ and $r_3$ of the form $f_i(z)g_i(x,y)$. Let $V$ be the space of
$h:A\to F$ with $\sum_xh(x)f_i(x)=0$ for the $r_1$ functions $f_i$ of the first kind. It is cut out by $r_1$ linear equations, so $\dim V\ge|A|-r_1$, and by Lemma~\ref{combinatorial:lem:support} it contains $h$ with
$|\operatorname{supp}h|\ge|A|-r_1$. Contract $T$ with $h$ in the variable $x$:
\[
U(y,z)=\sum_{x\in A}h(x)T(x,y,z)=\sum_{a\in A}c_ah(a)\,\delta_a(y)\delta_a(z).
\]
As a matrix indexed by $(y,z)$, $U$ is diagonal with $|\operatorname{supp}h|$ non-zero entries, so its rank is at least $|A|-r_1$. On the other side, contracting the slices: those of the first kind give
$\bigl(\sum_xh(x)f_i(x)\bigr)g_i(y,z)=0$; those of the second kind give $f_i(y)\sum_xh(x)g_i(x,z)$, and those of the third kind $f_i(z)\sum_xh(x)g_i(x,y)$, which are matrices of rank at most $1$. So
$\operatorname{rank}U\le r_2+r_3<|A|-r_1$, a contradiction.
\end{proof}

Let $M_n$ be the number of $a\in\{0,1,2\}^n$ with $a_1+\dots+a_n\le\frac{2n}3$.

\begin{theorem}[Ellenberg and Gijswijt, 2016]\label{combinatorial:thm:capset}
Every cap set $A\subseteq\mathbb F_3^n$ has $|A|\le3M_n$.
\end{theorem}
\begin{proof}
Over $F=\mathbb F_3$, let $T(x,y,z)=1$ if $x+y+z=0$ and $0$ otherwise, for $x,y,z\in A$. If $x+y+z=0$ with $x,y,z\in A$, then $x=y=z$: three distinct points would form a line, and if two of them are equal,
say $x=y$, then $z=-2x=x$. So $T(x,y,z)=\sum_{a\in A}\delta_a(x)\delta_a(y)\delta_a(z)$, and by Lemma~\ref{combinatorial:lem:diag} its slice rank is $|A|$.

Now bound the slice rank from above. For $t\in\mathbb F_3$, $1-t^2$ is $1$ if $t=0$ and $0$ otherwise. Hence
\[
T(x,y,z)=\prod_{i=1}^n\bigl(1-(x_i+y_i+z_i)^2\bigr).
\]
Expanding the product gives a combination of monomials $x^ay^bz^c=\prod_ix_i^{a_i}y_i^{b_i}z_i^{c_i}$ with all exponents in $\{0,1,2\}$ and total degree $|a|+|b|+|c|\le2n$, where $|a|=\sum_ia_i$. In each
monomial, one of $|a|,|b|,|c|$ is at most $\frac{2n}3$. Group the monomials with $|a|\le\frac{2n}3$ by their $x$-part: they contribute $\sum_ax^ag_a(y,z)$, a sum of at most $M_n$ slices of the first kind. Group the
remaining monomials with $|b|\le\frac{2n}3$ by their $y$-part, and the rest (which have $|c|\le\frac{2n}3$) by their $z$-part. So the slice rank is at most $3M_n$.
\end{proof}

\begin{proposition}\label{combinatorial:prop:mn}
$M_n\le\bigl(\frac{1+x+x^2}{x^{2/3}}\bigr)^n$ for every $0<x\le1$. With $x=0.6$, $M_n\le2.7552^n$, so every cap set in $\mathbb F_3^n$ has at most $3\cdot2.7552^n$ points.
\end{proposition}
\begin{proof}
For $0<x\le1$, each $a$ with $|a|\le\frac{2n}3$ has $x^{|a|-2n/3}\ge1$, so
\[
M_n\le\sum_{a\in\{0,1,2\}^n}x^{|a|-2n/3}=x^{-2n/3}\prod_{i=1}^n\sum_{a_i=0}^2x^{a_i}=\Bigl(\frac{1+x+x^2}{x^{2/3}}\Bigr)^n.
\]
At $x=0.6$, $\frac{1.96}{0.6^{2/3}}=\frac{1.96}{0.71138\ldots}=2.75521\ldots$
\end{proof}

This is the exponential trick behind Chernoff's bound (Volume~4, Chapter~4): the fraction $\frac{M_n}{3^n}$ is the probability that a sum of $n$ uniform digits in $\{0,1,2\}$ falls below $\frac23$ of its mean $n$,
which is exponentially small. The optimal $x=\frac{\sqrt{33}-1}8\approx0.5931$ gives $2.7551$, and the method cannot do much better: the same argument bounds the more general ``tri-coloured sum-free sets'', and for those Kleinberg, Sawin and Speyer (with Norin and Pebody) showed that
$2.7551^n$ is the true growth, up to smaller factors.
Whether the truth is near $2.22^n$, near $2.755^n$, or in between, nobody knows.

\begin{think}
The proof never used that $A$ has no line except through the shape of $T$ on $A\times A\times A$. Which other problems ``no solution of $x+y+z=0$ with $x,y,z$ not all equal'' can it solve? Try sets with no
three-term arithmetic progression in $\mathbb F_p^n$, for a prime $p\ge5$.
\end{think}

\begin{beauty}
Clubs, wheat fields, points in space, Ramsey colourings, sums modulo a prime and SET cards: in each case the objects were turned into vectors, the condition into independence, and the count into a dimension.
The method never looks at the objects themselves, never searches, never estimates; it notices that they cannot all fit. Its greatest triumph came in 2016, when a few lines about polynomials of
degree $2$ in each variable settled a problem that decades of Fourier analysis had only nibbled at.
\end{beauty}

\section{Exercises}

\begin{exercise}[$\star$]\label{combinatorial:ex:eventown}
Check that the couples construction satisfies the Eventown rules. Then let the inhabitants be the $8$ points of $\mathbb F_2^3$ and the clubs the $14$ affine planes (translates of $2$-dimensional subspaces),
together with $\varnothing$ and $\mathbb F_2^3$. Show that this is an Eventown family of maximum size which does not come from couples.
\end{exercise}

\begin{exercise}[$\star$]\label{combinatorial:ex:nearpencil}
Check Fisher's inequality is attained with $\lambda=1$ by the \emph{near-pencil} $\{1,n\},\{2,n\},\dots,\{n-1,n\},\{1,\dots,n-1\}$, and by the seven lines of the Fano plane (Volume~9, Appendix~C) on $n=7$ points.
\end{exercise}

\begin{exercise}[$\star$]\label{combinatorial:ex:simplexdist}
In $\R^{n+1}$, consider the points $e_i+e_j$, $1\le i<j\le n+1$. Show that they lie in a hyperplane and form a two-distance set of $\frac{n(n+1)}2$ points in $\R^n$. What is it for $n=3$?
\end{exercise}

\begin{exercise}[$\star\star$]\label{combinatorial:ex:sdistance}
An \emph{$s$-distance set} in $\R^n$ is a set whose pairwise distances take at most $s$ values. Show that it has at most $\binom{n+2s}{n}$ points.
\end{exercise}

\begin{exercise}[$\star$]\label{combinatorial:ex:fwmodthree}
Let $A_1,\dots,A_m\subseteq[n]$ have sizes divisible by $3$, while no $|A_i\cap A_j|$, $i\ne j$, is divisible by $3$. Show that $m\le1+n+\binom n2$.
\end{exercise}

\begin{exercise}[$\star$]\label{combinatorial:ex:cdsharp}
Show that equality holds in the Cauchy--Davenport theorem when $A$ and $B$ are arithmetic progressions with the same difference and $|A|+|B|\le p+1$. Show that the theorem fails in $\Z/4\Z$.
\end{exercise}

\begin{exercise}[$\star$]\label{combinatorial:ex:egzsharp}
Show that $2n-2$ integers need not contain $n$ whose sum is divisible by $n$.
\end{exercise}

\begin{exercise}[$\star\star$]\label{combinatorial:ex:quadric}
Let $p$ be a prime. Show that every quadratic form $\sum_{i\le j}c_{ij}x_ix_j$ in $n\ge3$ variables over $\mathbb F_p$ has a non-trivial zero. Deduce that $-1$ is a sum of two squares modulo $p$.
\end{exercise}

\begin{exercise}[$\star\star$]\label{combinatorial:ex:capplane}
Show that a cap set in $\mathbb F_3^2$ has at most $4$ points.
\end{exercise}

\begin{exercise}[$\star\star\star$]\label{combinatorial:ex:dvir}
\emph{Dvir's Kakeya theorem.} A set $K\subseteq\mathbb F_q^n$ is a \emph{Kakeya set} if it contains a line in every direction: for every $d\ne0$ there is $x$ with $\{x+td:t\in\mathbb F_q\}\subseteq K$. Show that
$|K|\ge\binom{q+n-1}n$. (Use a non-zero polynomial of degree at most $q-1$ vanishing on $K$, and look at its restriction to lines. Lemma~\ref{combinatorial:lem:grid} holds over any field.)
\end{exercise}

\section{Hints}

\begin{hint}{combinatorial:ex:eventown}
Two distinct affine planes of $\mathbb F_2^3$ are either parallel or meet in an affine line. How many $2$-dimensional subspaces does $\mathbb F_2^3$ have?
\end{hint}
\begin{hint}{combinatorial:ex:nearpencil}
Count intersections directly. In the Fano plane, two lines meet in exactly one point.
\end{hint}
\begin{hint}{combinatorial:ex:simplexdist}
The coordinates of $e_i+e_j$ add up to $2$. Two such points share one index or none.
\end{hint}
\begin{hint}{combinatorial:ex:sdistance}
Use $f_i(x)=\prod_{d}(|x-a_i|^2-d^2)$, a polynomial of degree $2s$.
\end{hint}
\begin{hint}{combinatorial:ex:fwmodthree}
Theorem~\ref{combinatorial:thm:fwmod} with $p=3$ and $L=\{1,2\}$.
\end{hint}
\begin{hint}{combinatorial:ex:cdsharp}
$\{0,d,\dots,(k-1)d\}+\{0,d,\dots,(l-1)d\}$. In $\Z/4\Z$, try a subgroup.
\end{hint}
\begin{hint}{combinatorial:ex:egzsharp}
Use only the values $0$ and $1$.
\end{hint}
\begin{hint}{combinatorial:ex:quadric}
Chevalley--Warning with one polynomial of degree $2<n$. For the second part, use $x^2+y^2+z^2$.
\end{hint}
\begin{hint}{combinatorial:ex:capplane}
In a cap of $5$ points, look at the four lines through one of its points, then at the four lines through a point outside it.
\end{hint}
\begin{hint}{combinatorial:ex:dvir}
The polynomials of degree at most $q-1$ in $n$ variables form a space of dimension $\binom{q-1+n}n$. On the line $x+td$, the restriction is a polynomial in $t$ of degree at most $q-1$, whose coefficient
of $t^{\deg f}$ is the top-degree part of $f$ evaluated at $d$.
\end{hint}

\section{Solutions}

\begin{solution}{combinatorial:ex:eventown}
\idea Unions of couples meet in unions of couples. For the planes, the intersection of two affine subspaces is empty or an affine subspace, whose size is a power of $2$.
\writeup Two unions of couples intersect in a union of couples, of even size. For the second family: a $2$-dimensional subspace of $\mathbb F_2^3$ is the kernel of a non-zero linear form, and there are $7$ of them, each with
$2$ cosets: $14$ planes of $4$ points. Two distinct planes are parallel (disjoint) or their directions differ, and then they meet in a translate of the intersection of the two directions, a line of $2$ points. With
$\varnothing$ and the whole space, which meet everything in $0$, $4$ or $8$ points, all sizes and intersections are even, and there are $16=2^{8/2}$ clubs: the maximum allowed by Theorem~\ref{combinatorial:thm:eventown}. No club has $2$
members, while the couples family contains the couples themselves.
\end{solution}

\begin{solution}{combinatorial:ex:nearpencil}
\idea Count the common points pair by pair.
\writeup $\{i,n\}\cap\{j,n\}=\{n\}$ for $i\ne j$, and $\{i,n\}\cap\{1,\dots,n-1\}=\{i\}$: every intersection has size $1$, and there are $n$ sets. In the Fano plane any two of the seven lines
$124,135,236,167,257,347,456$ meet in exactly one point (two points determine their line, so two lines share at most one point, and the $21$ pairs can be checked one by one: $124\cap456=\{4\}$, $124\cap236=\{2\}$, $167\cap257=\{7\}$, and so on). So $m=n=7$. The incidence calculation shows
that every projective plane attains the bound, and de Bruijn and Erd\H os (1948) proved that for $\lambda=1$ the near-pencils and the projective planes are the only families with $m=n$.
\end{solution}

\begin{solution}{combinatorial:ex:simplexdist}
\idea Points whose coordinates have a fixed sum live in a hyperplane; distances depend only on how many indices are shared.
\writeup All the points satisfy $\sum_kx_k=2$, a hyperplane of dimension $n$, which is isometric to $\R^n$. For $\{i,j\}\ne\{k,l\}$, $|e_i+e_j-e_k-e_l|^2$ is $2$ if the pairs share an index and $4$ if not: two distances,
$\binom{n+1}2$ points. For $n=3$ the six points $e_i+e_j$ of $\R^4$ form a regular octahedron: each point is at distance $\sqrt2$ from the four points sharing an index and at distance $2$ from the opposite one.
\end{solution}

\begin{solution}{combinatorial:ex:sdistance}
\idea Exactly as for two distances, with one factor per distance.
\writeup Let $a_1,\dots,a_m$ have distances in $\{d_1,\dots,d_s\}$ and $f_i(x)=\prod_{k=1}^s(|x-a_i|^2-d_k^2)$. Then $f_i(a_j)=0$ for $j\ne i$ and $f_i(a_i)=\prod_kd_k^2\ne0$, so the $f_i$ are independent. They are
polynomials of degree at most $2s$ in $n$ variables, and the monomials of degree at most $2s$ in $n$ variables are in bijection with the monomials of degree exactly $2s$ in $n+1$ variables (add a variable $x_0$
to complete the degree), counted by stars and bars: $\binom{2s+n}{n}$. For $s=2$ this is $\binom{n+4}4$, much weaker than Theorem~\ref{combinatorial:thm:twodist}, which used the special shape of $|x-a|^2$.
\end{solution}

\begin{solution}{combinatorial:ex:fwmodthree}
\idea Exactly the shape of the modular Frankl--Wilson theorem.
\writeup Take $p=3$ and $L=\{1,2\}$, so $s=2$. Sizes are $\equiv0\notin L$ and intersections are $\equiv1$ or $2$, in $L$. Theorem~\ref{combinatorial:thm:fwmod} gives $m\le\binom n0+\binom n1+\binom n2$.
\end{solution}

\begin{solution}{combinatorial:ex:cdsharp}
\idea Sums of two progressions with the same difference form a progression.
\writeup With $A=\{a,a+d,\dots,a+(k-1)d\}$ and $B=\{b,\dots,b+(l-1)d\}$, $d\ne0$, $A+B=\{a+b+jd:0\le j\le k+l-2\}$, which has $k+l-1$ elements when $k+l-1\le p$, since the multiples $jd$, $0\le j<p$, are distinct. In
$\Z/4\Z$, $A=B=\{0,2\}$ gives $A+B=\{0,2\}$, of size $2<\min(4,3)$: the subgroup $\{0,2\}$ is closed under addition, which cannot happen in $\Z/p\Z$.
\end{solution}

\begin{solution}{combinatorial:ex:egzsharp}
\idea If all values are $0$ or $1$, a sum of $n$ of them is divisible by $n$ only when they are all equal.
\writeup Take $n-1$ zeros and $n-1$ ones. A choice of $n$ of them contains $k$ ones with $1\le k\le n-1$ (there are not $n$ zeros or $n$ ones), and its sum $k$ is not divisible by $n$.
\end{solution}

\begin{solution}{combinatorial:ex:quadric}
\idea One polynomial of degree $2$ in at least $3$ variables, vanishing at $0$: exactly Chevalley--Warning's setting.
\writeup The form $Q$ has degree at most $2<3\le n$ and $Q(0)=0$, so by Theorem~\ref{combinatorial:thm:cw} it has a non-zero zero. For $Q=x^2+y^2+z^2$, take a zero $(x,y,z)\ne0$. If $z\ne0$, then $-1=(x/z)^2+(y/z)^2$. If $z=0$,
then $x^2+y^2=0$ with, say, $x\ne0$, and $-1=(y/x)^2+0^2$. (For $p=2$, $-1=1=1^2+0^2$.) This is the starting point of Lagrange's four-square theorem (Volume~2, Chapter~7), which also underpins the
Bruck--Ryser theorem.
\end{solution}

\begin{solution}{combinatorial:ex:capplane}
\idea A cap meets each line in at most $2$ points; with $5$ points the lines through a cap point are all ``full'', and then a parity count on the lines through a point outside the cap fails.
\writeup Suppose $A$ is a cap with $|A|=5$. The plane has $4$ directions, and through each point there is one line in each direction; the $4$ lines through a point partition the other $8$ points. Fix $P\in A$. Each line
through $P$ contains at most one other point of $A$, and the four other points of $A$ lie on these four lines, so each line through $P$ contains exactly one other point of $A$. As $P$ was arbitrary, every line
meeting $A$ meets it in exactly $2$ points. Now fix $Q\notin A$ (there are $4$ such points). The four lines through $Q$ partition $A$, and each meets $A$ in $0$ or $2$ points: so $|A|$ is even, a contradiction.
A cap of size $4$ exists (Figure~\ref{combinatorial:fig:caps}).
\end{solution}

\begin{solution}{combinatorial:ex:dvir}
\idea If $K$ is small, a low-degree polynomial vanishes on it. On each line of $K$ it vanishes at $q$ points, more than its degree, so it vanishes on the whole line, and its top part vanishes on every direction: too many zeros.
\writeup Suppose $|K|<\binom{q+n-1}n$. The space of polynomials of degree at most $q-1$ in $n$ variables has dimension $\binom{q-1+n}n>|K|$, and vanishing on $K$ is $|K|$ linear conditions, so some non-zero $f$ with
$\deg f=e\le q-1$ vanishes on $K$. If $e=0$, $f$ is a non-zero constant, which cannot vanish on $K\ne\varnothing$. So $e\ge1$; write $f_e$ for the homogeneous part of degree $e$. For a direction $d\ne0$, take $x$ with
the line $x+td$ in $K$; then $g(t)=f(x+td)$ is a polynomial of degree at most $e\le q-1$ in $t$, vanishing at all $q$ values of $t$, so $g=0$. Its coefficient of $t^e$ is $f_e(d)$, so $f_e(d)=0$. Also $f_e(0)=0$ as
$e\ge1$. So $f_e$ vanishes on $\mathbb F_q^n$, with $\deg_{x_i}f_e\le e<q=|\mathbb F_q|$, and Lemma~\ref{combinatorial:lem:grid} gives $f_e=0$, contradicting $\deg f=e$. Since $\binom{q+n-1}n\ge\frac{q^n}{n!}$, a Kakeya set contains a
fixed fraction of the space. In 2008 this two-paragraph argument settled the finite field version of a famous problem of harmonic analysis, which was thought to be as hard as the real one.
\end{solution}

\section*{Sources and further reading}

\begin{itemize}
\item \emph{The stories.} The epigraph is from E.~Artin, \emph{Geometric Algebra}, Interscience, 1957, Chapter~I. Fisher's years at Rothamsted are told in J.~Fisher Box, \emph{R.~A.~Fisher, the Life of a
Scientist}, Wiley, 1978. The planar two-distance bound is L.~M.~Kelly's solution of Erd\H os' problem E735, American Mathematical Monthly 54 (1947), 227. The origin of SET is told by its publisher, Set
Enterprises, and in L.~McMahon, G.~Gordon, H.~Gordon and R.~Gordon, \emph{The Joy of SET}, Princeton University Press, 2016, which also covers caps and the Ellenberg--Gijswijt bound.
\item L.~Babai and P.~Frankl, \emph{Linear Algebra Methods in Combinatorics}, preliminary version, University of Chicago, 1992 (2nd ed.\ 2020): the source of Oddtown and Eventown as named here, and of much of this
chapter; J.~Matou\v{s}ek, \emph{Thirty-three Miniatures}, AMS, 2010, short proofs in the same spirit; E.~R.~Berlekamp, \emph{On subsets with intersections of even cardinality}, Canadian Mathematical Bulletin 12
(1969), 471--474.
\item R.~A.~Fisher, \emph{An examination of the different possible solutions of a problem in incomplete blocks}, Annals of Eugenics 10 (1940), 52--75; R.~C.~Bose, \emph{A note on Fisher's inequality for balanced
incomplete block designs}, Annals of Mathematical Statistics 20 (1949), 619--620; N.~G.~de Bruijn and P.~Erd\H os, \emph{On a combinatorial problem}, Indagationes Mathematicae 10 (1948), 421--423.
\item D.~G.~Larman, C.~A.~Rogers and J.~J.~Seidel, \emph{On two-distance sets in Euclidean space}, Bulletin of the London Mathematical Society 9 (1977), 261--267; A.~Blokhuis, \emph{Few-distance sets}, CWI Tract
7, 1984.
\item P.~Frankl and R.~M.~Wilson, \emph{Intersection theorems with geometric consequences}, Combinatorica 1 (1981), 357--368; J.~Kahn and G.~Kalai, \emph{A counterexample to Borsuk's conjecture}, Bulletin of the AMS 29
(1993), 60--62, two pages; E.~Chattopadhyay and D.~Zuckerman, \emph{Explicit two-source extractors and resilient functions}, Annals of Mathematics 189 (2019), 653--705.
\item N.~Alon, \emph{Combinatorial Nullstellensatz}, Combinatorics, Probability and Computing 8 (1999), 7--29, with many more applications; P.~Erd\H os, A.~Ginzburg and A.~Ziv, \emph{Theorem in the additive number
theory}, Bulletin of the Research Council of Israel 10F (1961), 41--43; Z.~Dvir, \emph{On the size of Kakeya sets in finite fields}, Journal of the AMS 22 (2009), 1093--1097.
\item \emph{Research papers to read after this chapter.} E.~Croot, V.~Lev and P.~P.~Pach, \emph{Progression-free sets in $\Z_4^n$ are exponentially small}, Annals of Mathematics 185 (2017), 331--337, and J.~S.~Ellenberg
and D.~Gijswijt, \emph{On large subsets of $\mathbb F_q^n$ with no three-term arithmetic progression}, the same volume, 339--343. Read Ellenberg--Gijswijt first: its Lemma~1 is a version of our slice rank bound
written with a matrix of polynomials (the ``polynomial method'' of Croot--Lev--Pach), and its main proof is our Theorem~\ref{combinatorial:thm:capset}, with $\mathbb F_q$ for $\mathbb F_3$ and the general count of monomials. Then
T.~Tao, \emph{A symmetric formulation of the Croot--Lev--Pach--Ellenberg--Gijswijt capset bound}, blog post, 18 May 2016, which introduced slice rank, and R.~Kleinberg, W.~Sawin and D.~Speyer, \emph{The growth of
tri-colored sum-free sets}, Discrete Analysis 2018:12, on the limits of the method. For the history and the lower bounds: R.~Meshulam, \emph{On subsets of finite abelian groups with no 3-term arithmetic
progressions}, Journal of Combinatorial Theory A 71 (1995), 168--172; Y.~Edel, S.~Ferret, I.~Landjev and L.~Storme, \emph{The classification of the largest caps in AG(5,3)}, Journal of Combinatorial Theory A 99
(2002), 95--110; B.~Romera-Paredes et al., \emph{Mathematical discoveries from program search with large language models}, Nature 625 (2024), 468--475, which found the current lower bound $2.2202^n$.
\end{itemize}
